\section{Introduction}
\label{sec:intro}

Robots and wearable augmented-reality systems often need a map before data
collection is complete. Reconstructing scene appearance supports visualization
from new viewpoints, while reconstructing geometry supports occlusion and
spatial interaction. For these applications, the map must be refined as new
observations arrive. High reconstruction quality achieved only after an
offline refinement stage does not meet this online requirement.

A map represented by explicit, colored Gaussian primitives can be rendered
efficiently and refined through image-based optimization. This representation,
3D Gaussian Splatting (3DGS), combines high-quality view synthesis with fast
rendering~\cite{kerbl20233dgs}. Recent SLAM systems have extended it to map
construction during camera tracking~\cite{matsuki2024monogs,splatam2024,
huang2024photoslam,rtgslam2024}. However, fast rendering does not imply fast
map refinement. A system can display its current map at a high frame rate
while recently observed regions remain poorly reconstructed. Online mapping
must also share computation with tracking, limiting the number of map
optimization steps, or \emph{mapping updates}, that can be completed as the
stream arrives. The challenge is therefore to obtain a useful map from these
limited updates, not merely to render it quickly.

Existing methods improve the use of online computation in several ways.
Photo-SLAM progressively refines scene appearance using geometry-based
densification and Gaussian-pyramid learning~\cite{huang2024photoslam}.
CaRtGS accelerates backpropagation and addresses uneven keyframe optimization
by assigning additional iterations according to rendering loss~\cite{cartgs}.
These advances establish that both the cost of each update and its allocation
matter. A further question concerns the observations used for mapping:
\emph{when should additional images enter optimization, and how should updates
be distributed across the resulting growing set?}

Images selected for camera tracking, called \emph{keyframes}, need not provide
all the observations useful for appearance reconstruction. Frames between
keyframes can supply additional viewpoints and texture evidence. We refer to
the images used for map optimization as \emph{training views}. Adding more
training views provides more supervision, but also distributes the available
mapping updates over a larger set. Moreover, uniform selection from the
current set does not equalize accumulated training: views added early remain
eligible for longer, whereas later views have fewer opportunities to be
selected. CaRtGS addresses uneven optimization within a keyframe
pool~\cite{cartgs}; here, we consider both expansion beyond tracking keyframes
and update allocation within that expanded training-view set. These decisions
must be made without knowing the remaining stream length or future
optimization capacity.

Better appearance alone is not sufficient for a reliable map. Gaussians at
incorrect depths can reproduce observed colors, leaving spurious structure
between the camera and a real surface~\cite{deng2022dsnerf,stablegs}. Depth
and normal constraints help resolve this ambiguity~\cite{huang20242dgs,
chen2024pgsr}. However, matching a single average rendered depth can still
leave incorrect contributions in front of and behind the surface. This
motivates directly constraining where the map places opacity along a viewing
ray. In an online system, such constraints must use depth evidence already
available at the current time.

Our aim is to improve map appearance and geometry under the limited mapping
updates available during online operation. We build on the visual--inertial
tracking and Gaussian mapping pipeline of VIGS-SLAM~\cite{vigsslam}.
First, we expand the training-view set beyond tracking keyframes only as
mapping updates are completed. This rule, \emph{Optimization-Guided View
Growth}, ties the addition of observations to completed optimization work.
Second, we give higher selection probabilities to views that have received
fewer updates. We retain randomness so that this preference does not reduce
selection to a deterministic least-count rule. Our \emph{Entropy-Regularized
View Sampling} expresses this trade-off through a count-dispersion objective
and an entropy term. Sampling without replacement further prevents repeated
selection of the same view within a group of updates. Finally, we use
available depth evidence to distinguish empty space in front of observed
surfaces from the surfaces themselves. We formulate corresponding constraints
on Gaussian opacity along viewing rays, rather than relying on color
agreement alone.

We compare rendering quality with VIGS-SLAM on RPNG, UTMM, and Aria.
Separate ablations examine the use of frames between keyframes at a fixed
update budget and compare view-selection policies at several update budgets.
These experiments assess how the choice and allocation of training views
affect the quality obtained from limited mapping updates.

Our contributions are:
\begin{itemize}
    \item Optimization-Guided View Growth, which expands the training-view
    set beyond tracking keyframes according to completed mapping updates.
    \item Entropy-Regularized View Sampling, which favors less frequently
    selected training views while retaining randomized selection and avoiding
    repetition within each group of updates.
    \item Depth-supported constraints on Gaussian opacity along viewing
    rays, formulated to distinguish observed free space from surface evidence
    using only observations available at the current time.
\end{itemize}

% AUTHOR NOTE: Sec. 3.3 still contains Carve, Hit, and Carve+Hit candidates.
% Specify the selected geometry objective here after that choice is finalized.
% Do not claim faster geometric convergence or time-resolved quality gains
% until the corresponding measurements are reported in the experiments.
